\section{Zaključak}

Ovaj rad predstavio je sistematičan pregled, implementaciju i komparativnu analizu metoda interpretabilnosti primenjenih na Vision Transformer (ViT) modele. Evaluacija je sprovedena na javno dostupnom ViT-Base/16-224 modelu i validacionom skupu ImageNet, korišćenjem kvantitativnih metrika (MoRF, LeRF, \textit{Pointing Game}, Spirmanova korelacija), kao i testova robustnosti specifičnih za Transformer arhitekture.

Eksperimentalni rezultati ukazuju na nekoliko ključnih zaključaka. Među metodama sa pristupom unutrašnjoj strukturi modela, GMAR se pokazao kao najbolji kompromis između kvaliteta mapa i računarske efikasnosti, uz dodatnu prednost izrazite robustnosti na perturbacije ulaza. Transformer Attribution postiže generalno najkvalitetnije mape i jedina je metoda koja dosledno ispoljava klasnu specifičnost, ali je za skoro red veličine sporija od GMAR-a. Attention Rollout se, uprkos svojoj brzini, pokazao kao nepouzdana metoda čiji rezultati konzistentno zaostaju za ostalima — što potvrđuje teorijsku pretpostavku da sirove mape pažnje nisu dovoljan signal za interpretabilnost. LRP se pokazao kao metoda izrazito osetljiva na lokalne perturbacije ulaza, što je u suprotnosti sa njenim teorijskim svojstvima konzervacije relevantnosti. Gradijentne metode (Vanilla Gradient, Input $\times$ Gradient, GradCAM) pokazuju umerene i međusobno slične rezultate, pri čemu GradCAM neznatno prednjači.

Među metodama crne kutije, RISE se pokazao kao ubedljivo superioran — postiže najbolje rezultate na MoRF/LeRF i \textit{Pointing Game} metrici od svih evaluiranih metoda, uključujući i metode sa pristupom modelu. Međutim, vreme izvršavanja od skoro četrdeset sekundi po slici čini ga nepraktičnim za primene koje zahtevaju obradu u realnom vremenu. Kernel SHAP se, uprkos formalnoj aksiomatskoj osnovi, pokazao kao neprikladna metoda za interpretabilnost slika — reprezentacija na nivou pečeva i inherentna pretpostavka aditivnosti doprinosa značajno ograničavaju njegovu efektivnost u vizuelnom domenu.

Generalno niske vrednosti Spirmanove korelacije između svih parova metoda ukazuju na to da različiti pristupi mere fundamentalno različite aspekte ponašanja modela, što otežava donošenje jednoznačnih zaključaka oslanjanjem na samo jednu metodu i naglašava važnost primene komplementarnih pristupa.

Doprinosi ovog rada ogledaju se u sistematičnom poređenju devet metoda interpretabilnosti u identičnim eksperimentalnim uslovima, što nije prisutno u postojećoj literaturi koja metode uglavnom evaluira izolovano. Pored toga, predloženi testovi robustnosti specifični za ViT arhitekturu — test globalnog mešanja pečeva i test nasumičnog peča — predstavljaju nov metodološki okvir koji uvažava specifičnosti mehanizma pažnje.

Rad ima i određena ograničenja. Evaluacija je sprovedena na jednom modelu i jednom skupu podataka, te generalizacija zaključaka na druge arhitekture i domene zahteva dodatna istraživanja. Pored toga, MoRF i LeRF metrike pate od problema narušavanja prirodne distribucije podataka usled maskiranja pečeva, što može uticati na objektivnost dobijenih rezultata.

U pogledu budućih istraživanja, ovaj rad se fokusirao na naknadnu interpretabilnost već istreniranog modela. Noviji pravci teže ka integraciji interpretabilnosti direktno u proces treniranja, čime se dobijaju modeli koji su inherentno interpretabilni bez potrebe za naknadnom analizom \cite{chen2019looks}. Pored toga, dok ovaj rad prati tok informacija od izlaza ka ulazu, postoje kauzalni pristupi koji prate suprotan smer — od ulaza ka izlazu — ispitujući direktno kako promene ulaza uzrokuju promene u odluci modela, što predstavlja konceptualno drugačiji i komplementaran pristup \cite{pearl2009causality}.